\chapter{Sự Thật Phũ Phàng}
\markboth{Sự Thật Phũ Phàng}{Sự Thật Phũ Phàng}

\begin{center}
	\textit{\color{bookprimary}``Ngoài đời không thiếu người nói hay. Ngoài đời chỉ thiếu người làm ra giá trị đủ rõ để người khác không thể làm ngơ.''}

	\textit{\color{bookprimary}``The world is not short of people who speak well. It is short of people who create value so clearly that others cannot ignore it.''}
\end{center}

\ngancach

\begin{tcolorbox}[
	colback=bookprimary!6,
	colframe=bookprimary!35,
	arc=5pt,
	title={\bfseries Góc Suy Nghĩ},
	fonttitle=\bfseries\color{bookprimary}
]
Nhiều bạn trẻ lớn lên với một ảo tưởng ngọt ngào: rồi sẽ có ai đó nhìn ra tiềm năng của mình, kéo mình lên, chỉ mình cách làm giàu, rồi mở cửa cho mình đi vào.
\textit{(Many young people grow up with a sweet illusion: that someday someone will notice their potential, pull them upward, teach them how to get rich, and open the door for them.)}

Sự thật phũ phàng hơn nhiều.
\textit{(The truth is much harsher.)}

Không ai rảnh đi làm giàu cho bạn cả.
\textit{(No one is out there with free time to make you rich.)}

Người ta chỉ hợp tác, chỉ chỉ đường, chỉ trao cơ hội khi họ thấy ở bạn một giá trị đủ thật để việc giúp bạn cũng có ý nghĩa với họ.
\textit{(People only collaborate, guide, or offer opportunities when they see enough real value in you for helping you to make sense for them as well.)}
\end{tcolorbox}

\section{Không Ai Rảnh Đi Làm Giàu Cho Bạn}
\begin{truyen}{Không Ai Rảnh Đi Làm Giàu Cho Bạn}{No One Is Busy Making You Rich}
\chuhoa{C}{ó bạn hỏi thầy rất thật}: ``Sao em thấy trên mạng ai cũng nói có bí kíp, có mentor, có người dẫn đường, mà ngoài đời em chẳng thấy ai thật lòng kéo em lên hết?''
\textit{(A student asks very honestly: ``Why does social media make it look like everyone has a secret formula, a mentor, a guide, but in real life I do not see anyone sincerely pulling me upward?'')}

Thầy đáp luôn: ``Vì ngoài đời, ai cũng đang bận kéo chính họ trước đã.''
\textit{(The teacher answers immediately: ``Because in real life, everyone is busy pulling themselves upward first.'')}

Bạn ấy ngơ ra.
\textit{(The student goes blank for a second.)}

\teacherpair{Em đừng tưởng thế giới là một nơi đầy những người rảnh rỗi đi tìm người có tiềm năng để nâng đỡ vô điều kiện. Người lớn ngoài kia còn phải lo tiền, lo việc, lo gia đình, lo giữ chỗ đứng của chính họ. Nếu họ giúp em, thường là vì họ thấy em đáng để giúp, hoặc giúp em cũng có ích cho công việc của họ.}{``Do not imagine the world as a place full of people with spare time searching for potential youngsters to support unconditionally. Adults out there are busy with money, work, family, and protecting their own footing. If they help you, it is usually because they see that you are worth helping, or because helping you is also useful to their work.''}

Thầy nói tiếp: ``Câu này nghe lạnh nhưng đúng: người chưa tự kéo nổi bản thân lên khỏi sự lười, sự mơ hồ và thói quen nói hay hơn làm thì rất khó gặp được người muốn kéo tiếp.''
\textit{(The teacher continues: ``This sounds cold, but it is true: people who cannot even pull themselves out of laziness, vagueness, and the habit of speaking better than they act are very unlikely to meet someone willing to pull them further.'')}
\end{truyen}

\section{Không Ai Giúp Bạn Mãi Nếu Bạn Không Đáng Để Giúp}
\begin{truyen}{Không Ai Giúp Bạn Mãi Nếu Bạn Không Đáng Để Giúp}{Help Follows Value}
\chuhoa{M}{ột bạn than}: ``Ban đầu ai cũng hứa hỗ trợ em. Nhưng sau một thời gian họ im luôn. Em thấy người lớn toàn nói hay thôi thầy.''
\textit{(A student complains: ``At first everyone promised to support me. But after a while they all went silent. Adults really only talk nicely, teacher.'')}

Thầy hỏi lại: ``Trong thời gian người ta đang hỗ trợ, em đã làm được gì để họ thấy công của họ không đổ sông đổ biển?''
\textit{(The teacher asks back: ``While they were supporting you, what did you actually do to show them their effort was not being wasted?'')}

Bạn ấy im lặng.
\textit{(The student falls silent.)}

\teacherpair{Giúp đỡ không phải cái vòi nước mở mãi. Người tử tế cũng mệt nếu họ cứ rót thời gian, lời khuyên và cơ hội vào một người không chịu lớn lên. Em muốn được giúp dài lâu thì em phải cho người ta thấy em biết nhận, biết làm, biết tiến bộ, và biết tự đứng dần lên.}{``Help is not a tap that runs forever. Even kind people get tired if they keep pouring time, advice, and opportunities into someone who refuses to grow. If you want long-term support, you must show people that you know how to receive, how to work, how to improve, and how to stand increasingly on your own.''}

Rồi thầy chốt: ``Nhiều người không thiếu người giúp. Họ thiếu tư cách để giữ sự giúp đỡ đó ở lại.''
\textit{(Then the teacher closes with: ``Many people do not lack helpers. They lack the character needed to keep that help from walking away.'')}
\end{truyen}

\section{Không Ai Dạy Bạn Làm Giàu Miễn Phí Rồi Ngồi Nhìn Bạn Vượt Mặt Họ}
\begin{truyen}{Không Ai Dạy Bạn Làm Giàu Miễn Phí Rồi Ngồi Nhìn Bạn Vượt Mặt Họ}{No One Hands Out Wealth for Free}
\chuhoa{C}{ó bạn mê xem clip làm giàu ngắn}: ``Em nghĩ chỉ cần gặp đúng người chỉ đường là đời em bẻ lái được luôn.''
\textit{(A student who loves short wealth-content clips says: ``I think if I just meet the right person to show me the way, my whole life can turn immediately.'')}

Thầy cười nhạt: ``Người chỉ đường thật thường không rảnh đứng giữa đường gào tên em.''
\textit{(The teacher smiles dryly: ``People who truly know the way are usually not standing in the middle of the road shouting your name.'')}

Bạn ấy cãi: ``Nhưng nhiều người thành công vẫn dạy miễn phí mà thầy.''
\textit{(The student pushes back: ``But many successful people still teach for free, teacher.'')}

\teacherpair{Có. Nhưng em phải hiểu họ đang cho em cái gì. Phần miễn phí thường chỉ là cánh cửa, là góc nhìn, là mồi suy nghĩ. Còn phần xương sống của việc làm ra tiền là năng lực, kỷ luật, sức bền, độ lì, khả năng chịu sai, và uy tín để người khác còn muốn quay lại làm ăn với em. Không ai gói mấy thứ đó thành quà rồi đặt vào tay em được.}{``Yes. But you need to understand what they are actually giving you. The free part is usually just a door, a perspective, a spark for thought. The backbone of making money is ability, discipline, stamina, toughness, tolerance for mistakes, and the credibility that makes people want to work with you again. No one can wrap those things as a gift and place them in your hands.''}

Thầy nói thêm: ``Nếu chỉ nghe vài câu truyền cảm hứng mà giàu được, thế giới này đã đầy triệu phú từ lâu rồi.''
\textit{(The teacher adds: ``If a few motivational lines were enough to make people rich, this world would have been overflowing with millionaires long ago.'')}
\end{truyen}

\section{Muốn Đổi Đời, Hãy Đổi Giá Trị Của Mình Trước}
\begin{truyen}{Muốn Đổi Đời, Hãy Đổi Giá Trị Của Mình Trước}{Change Your Value Before Asking for a Better Life}
\chuhoa{N}{hiều bạn nói rất lớn}: ``Em muốn đổi đời.''
\textit{(Many young people say loudly: ``I want to change my life.'')}

Thầy hỏi: ``Đổi bằng gì?''
\textit{(The teacher asks: ``Change it with what?'')}

Bạn ấy tưởng thầy hỏi khó.
\textit{(The student thinks it is a trick question.)}

Thầy giải thích: ``Đời chỉ đổi khi giá trị em mang ra trao đổi đổi. Em biết làm gì hơn trước, chịu được gì hơn trước, giữ chữ tín hơn trước, và tạo ra kết quả gì rõ hơn trước? Nếu mấy thứ đó không đổi mà em chỉ đổi khẩu hiệu, đổi ảnh nền, đổi thần tượng làm giàu, thì đời em vẫn đứng đó thôi.''
\textit{(The teacher explains: ``Life changes only when the value you bring into exchange changes. What can you do better than before? What can you endure better than before? How much better do you keep your word? What clearer results can you create? If those things do not change, and you only change slogans, wallpapers, or rich idols, then your life will remain exactly where it is.'')}

Rồi thầy nói một câu rất thô: ``Ngoài đời, ít ai quan tâm em khát vọng bao nhiêu. Người ta nhìn xem em giải quyết được việc gì.''
\textit{(Then the teacher says something very blunt: ``In real life, very few people care how ambitious you feel. They look at what problems you can solve.'')}
\end{truyen}

\begin{baihoc}
Muốn được kéo lên, trước hết phải bớt là gánh nặng.
\textit{(If you want to be pulled upward, first stop being a burden.)}

Muốn có người chỉ đường, trước hết phải đi được những bước cơ bản mà không cần ai bế.
\textit{(If you want someone to guide you, first prove that you can walk the basic steps without being carried.)}

Thế giới không nợ bạn sự giàu có chỉ vì bạn rất muốn nó.
\textit{(The world does not owe you wealth just because you want it badly.)}
\end{baihoc}

\begin{ghinhoanh}
\textbf{Short English to remember:}

No value, no leverage.

Ambition is cheap without proof.
\end{ghinhoanh}

\begin{tuhocnhanh}
1. Em đang chờ một người kéo mình lên, hay đang tự làm mình đáng để được kéo lên? \textit{(Are you waiting for someone to pull you up, or are you making yourself worth pulling up?)}

2. Nếu một người giỏi gặp em hôm nay, điều gì ở em khiến họ muốn giữ liên hệ? \textit{(If a capable person met you today, what in you would make them want to stay connected?)}

3. Em đang muốn giàu vì thích cảm giác được hơn người, hay vì thật sự có thứ giá trị muốn tạo ra? \textit{(Do you want wealth because you like the feeling of being above others, or because you truly want to create something valuable?)}
\end{tuhocnhanh}

\begin{tongketchuong}
Nói thẳng ra, đa số người trẻ không thiếu động lực miệng.
\textit{(To put it bluntly, most young people do not lack verbal motivation.)}

Các em thiếu năng lực đủ rõ để người khác thấy hợp tác với mình là một khoản đầu tư đáng giá.
\textit{(What they lack is the kind of clear ability that makes others see working with them as a worthwhile investment.)}

Sự thật phũ phàng không phải để làm nản.
\textit{(This harsh truth is not meant to discourage.)}

Nó để cắt bỏ ảo tưởng rằng chỉ cần muốn nhiều là sẽ được đời ưu tiên nhiều.
\textit{(It is meant to cut away the illusion that wanting something badly is enough to make life prioritize you.)}
\end{tongketchuong}

\begin{goigiaovien}
Nếu dùng chương này cho học sinh lớp 11, 12, nên giữ giọng thẳng nhưng không mạt sát.
\textit{(If this chapter is used with grade 11 or 12 students, keep the tone direct but not degrading.)}

Mục tiêu không phải dập tắt khát vọng làm giàu.
\textit{(The goal is not to crush the desire for wealth.)}

Mục tiêu là kéo khát vọng đó về gần với lao động thật, giá trị thật và năng lực thật.
\textit{(The goal is to bring that ambition closer to real work, real value, and real ability.)}

Kể chương này tốt nhất khi đi cùng câu hỏi rất cụ thể: hôm nay em đang làm gì để mình bớt là người chỉ biết muốn?
\textit{(This chapter works best when paired with one very concrete question: what are you doing today so that you become less of a person who only knows how to want?)}
\end{goigiaovien}

\ngancach